
\cite{Zhao2016} : Peng-Robinson-Stryjek-Vera EoS with modified Wong-Sandler mixing rule (binary interaction parameter composition linearly dependent)

- Calcula AAD de cada referência mostrando range de T e P

-REFS de outros trabalhos queue fizeram ELV CW com diferentes EoS

-The proposed EOS/Gex model has four adjustable temperature-dependent parameters for polar molecules. Once aij is fixed, the proposed model has four adjustable parameters for polar or associating fluid mixture per binary components: kij and kji for the WS mixing and combing rule, tij and tji for
the modified NRTL model.

-An excellent result was obtained when compared the modeling results with a large amount of the vapor-liquid equilibria experimental data (more than 1300
experimental data points located in a P-T region of 273e623 K and 0.1e200 MPa) for the CO 2 eH 2 O system.

-Therefore, for compositional reservoir simulation studies a reliable and efficient phase behavior model of the CO 2 eH 2 O-Hydrocarbon system is still needed to achieve more realistic simulation results.

- Panagiotopoulos and Reid [11] introduced a second interaction parameter in the vdW one-fluid mixing rule and achieved more successful results than does the original vdW mixing rule in some cases including the CO2-H2O system.

- Bjørner and Kontogeorgis [39] incorporated a quadrupolar contribution into the CPA EOS (resulting qCPA EOS) to improve the performance of CPA EOS on various CO 2 -related binary systems

- based on our analysis of reliability of the published VLE measurements on high pressures and high temperatures we have selected to use in this study the  experimental data reported by  Todheide and Franck [47] and Takenouchi and Kennedy [59], while ignoring the experimental data reported by Takenouchi and Kennedy [64].




\cite{Tsivintzelis2011}: CPA Eos. Both phase equilibria (vapor–liquid and liquid–liquid) and densities are considered for the mixtures involved. Different approaches for modeling pure CO2 and mixtures are compared. CO 2 is modeled as non self-associating fluid, or as self-associating component having two,
three and four association sites. It is shown that overall very good correlation is obtained for binary mixtures of CO 2 and water or alcohols when the solvation between CO 2 and the polar compound is explicitly accounted for (cross-association), whereas the model is less satisfactory when CO 2 is treated as self-associating compound. - 298 < T < 477 K

- phase equilibria of supercritical fluids including acid gases like CO 2 and H 2 S in mixtures with water, alcohols and glycols are of great importance in numerous processes in the oil and gas and chemical industries [2–4]. Water, typically found in petroleum reservoirs, or brine may cause reduction of the amount of gas available for mixing with hydrocarbons and this effect increases with pressure and the amount of the aqueous phase (while it decreases with salinity). In the chemical industry, supercritical fluid extraction is a separation process which may in the future replace, to some extent, environmentally less friendly processes such as distillation and
liquid–liquid extraction. One potential application is the recovery
of alcohols especially ethanol from aqueous solutions using CO 2
[5].

-The nature of CO 2 –water interactions has been studied both
from experimental and theoretical (mainly using ab initio calcu-
lations) point of view. Such interactions are due to dispersive
and dipole–quadrupole forces. Recently, a Lewis acid–base type of
interaction has been put in evidence [9]. In this, the carbon atom in
CO 2 molecule acts as electron acceptor, while the oxygen atom in
H 2 O molecule plays the role of an electron donor. On the other hand,
hydrogen bonding interactions between the CO 2 oxygen atoms and
the H 2 O protons can also be present, but are less stable than the
electron donor–acceptor interactions [9,10].

-An alternative approach to modeling CO 2 systems within the
cubic EoS framework is via the so-called EoS/G E mixing rules. While an “overall” description is often satisfactory with EoS/G E
models, they fail when a closer look at details of phase behav-
ior and extreme conditions is focused.

-otimo review sobre cubicas, SAFT e CPA ja usadas. Também calcula densidade da mistura an saturação e compare com EXP

- A common denominator in all the aforementioned investigations is that capturing the temperature-dependency of the solubilities involving water and acid gases is a difficult task, especially when this is attempted with temperature-independent parameters.



\cite{Diamontonis2012}:CO 2 is modeled as a non-associating or associating component within the Perturbed Chain-SAFT (PC-SAFT) and as a quadrupolar component within the truncated PC-Polar SAFT (tPC-PSAFT). It is shown that PC-SAFT with explicit account of H 2 O–CO 2 cross-association and tPC-PSAFT with explicit account of CO 2 quadrupolar interactions are the most accurate of the models examined. Saturated liquid mixture density data are accurately predicted by the two models.

- 298 < T < 533 K

- table 4 shows AAD reported for each work

-Modeling CO 2 as an associating component provides a better correlation of pure component data than modeling it as a quadrupolar fluid. For mixture calculations, explicit account of H 2 O–CO 2 solvation effects through an association scheme clearly improves VLE and LLE correlation. A temperature-independent binary interaction parameter is sufficient in all cases. Addition of temperature dependence in this parameter provides a marginal
improvement in the calculations.




\cite{Li2009}: We have recently proposed an accurate version of the cubic-plus-association (CPA) equation of state (EOS) for water-containing mixtures which combines the Peng-Robinson equation (PR) for the physical interactions and the thermodynamic perturbation theory for the hydrogen bonding of water molecules. Despite the significant improvement, the water composition in the nonaqueous phase is systematically underestimated for some systems where the nonwater species are methane and ethane at very highpressures, unsaturated hydrocarbons, CO 2 , and H 2 S. We attribute the deficiency to the neglect of the ‘‘cross association’’ between water and those nonwater molecules. In
this work, the accuracy is drastically improved by treating methane, ethane, unsaturated hydrocarbons, CO 2 and H 2 S as ‘‘pseudo-associating’’ components and describing the cross association with water in the framework of the perturbation theory.

- Similar to water molecules, each pseudo-associating molecule is assumed to possess four association sites belonging to two types (a i and b I ) with two sites for each type. there is neither cross association nor self association between pseudoassociating components, i.e., the ‘‘association strength’’ D a i b j ¼ 0. Second, the cross association is assumed to be symmetric between two sites of different types of water and pseudo-
associating component i, i.e., D ab i ¼ D a i b .

- 304 < T < 533 K



 \cite{Pappa2009}: BOA INTRO: The considered models include the Peng–Robinson equation of state coupled with the conventional van der Waals one fluid mixing rules or the universal mixing rules, and the cubic-plus-association equation of state that uses the Peng–Robinson equation of state in order to account for the usual attractive and repulsive forces and an extra association term to account for the strong hydrogen bonding interactions. The required model parameters are calibrated using experimental data up to 1500 bar for pressure, and up to 673 K for temperature. To improve the accuracy of the proposed models we consider two temperature ranges. Temperatures lower than 373 K are of interest to the geological CO 2 sequestration, while higher temperatures are of interest to fluid-inclusion studies.

-Blencoe et al. [46] and Blencoe [47] focused their discussion on higher temperatures (above 383 K) and pressures that are of interest to fluid-inclusion studies. Based on these works, the experimental data base shown in Table 3 was compiled in this study. USA TAKENOUSHI, E NÃO TODHEIDE PARA t ALTO

-Overall the CPA-PR was found to have the best performance.
In particular, the error in the CO 2 solubility in the aqueous phase
was calculated 5.9 \% and the H 2 O solubility in the CO 2 -rich phase
was calculated 9.6 \% for temperatures lower than 373 K. The cor-
responding errors for temperatures in the range 373–623 K were
calculated 7.1 and 10.3 \%, respectively.

- Usa kij e Iij (para a e b, respectivamente) sendo dependents de T e apresentando duas equações, uma para T < 373 K e outra para T > 373 K



\cite{Monteiro2020}: The CPA EoS applicability was demonstrated for the description of phase behavior and volumetric properties of CO 2 +water and
CO 2 +ethanol and pure compounds over an extensive range of
temperature and pressure. 4C association scheme was used for
water and carbon dioxide while 2B for ethanol. Pure component
parameters were retrieved from literature and binary interaction
parameters were fitted to experimental phase equilibrium data. For
CO 2 +water, a quadratic temperature dependency on the interac-
tion parameter was found necessary. In relation to the binary systems, CO 2 +water VLE data were adequately described by CPA EoS. The liquid phase composition was calculated with errors less than 7\%. The vapor composition
presented deviations mostly in the range of 15\%. low accuracy of pure component critical point and binary mixtures critical locus description





EOS CO2-BRINE


\cite{Olsen2021}:  eCPA + HV (NRTL) mixing rule: Compared different DTU eCPA parameterizations. A new parametrization is proposed, that is able to reproduce the density without the need of a volume translation, and without loss of accuracy of other properties (ACTUALLY IT DOES NOT DESCRIBE DIELECTRIC SATURATION)

-Some of them contain the Born term, some not, some are based on Debye-Hückel equation, others use MSA (mean spherical approximation). All versions of e-CPA employ different parameter estimation approaches

- The benefit of having a constant volume translation like
the Peneloux is that it will not have an effect on properties like
activity coefficients and phase equilibria [23]

-The association scheme of CO 2 within CPA has been parameterized
in several ways for the system water-CO 2 [28] . In this work CO 2 is
assumed to be non-associating and without considering any cross-
association between CO 2 and water.


\cite{Maribo2015}:
In this work, the Cubic Plus Association (CPA) Equation of State is extended to handle mixtures containing electrolytes by including the electrostatic contributions from the Debye–Huckel and Born terms using a self-consistent model for the static permittivity

- The main advantage of the nonprimitive modeling
approach is that it does not require a model for the static
permittivity. Still, as the nonprimitive approach requires the
inclusion of dipole-dipole interactions, it is incompatible
with most modern equations of state, such as the PC-SAFT
and CPA that do not explicitly account for the dipole-dipole
interactions. The PC-SAFT or the CPA EoS may still be extended to handle electrolytes using the primitive models,

-In addition to the ion-ion interactions, another important
characteristic of electrolyte mixtures is that ions tend to go
toward the phase with the highest static permittivity. A con-
sequence of this is that ions display a large negative value of
the Gibbs energy of hydration, and this effect may be mod-
eled using the Born term, 19 or from the ion-dipole interac-
tions in the nonprimitive MSA. 10 While some researchers
have included a Born model in the electrolyte EoS, 20 it has
been often omitted from published electrolyte EoS. A com-
mon argument for omitting the Born term is that it will not
provide a contribution to the activity coefficients—however,
this is only the case when the model for the static permittivity is assumed to be independent of the mixture composition

- a new theoretical model
was developed to extend the framework by Onsager, 24 Kirk-
wood, 25 and Fr€ olich 26 for modeling of the static permittivity
to hydrogen-bonding compounds and salts. The new model
uses a simplified geometrical configuration of the dipolar
molecules in a hydrogen-bonding network to describe the
short-range dipolar correlations, and relies on the Wertheim
association theory as used in the SAFT-based equations of
state to determine the probability of association between dif-
ferent molecules.

-As the Born term is an internal energy contribution, it was decided
to fit the Born radii of the ions to data for the enthalpy of
hydration at 25 C. One of the challenging tasks in dealing with modeling of
electrolyte mixtures is how to obtain ion-specific parameters. For simple binary mixtures, salt-specific parameters are
adequate, but when dealing with mixtures containing multi-
ple ionic species, the system becomes over-specified.

- The model can be parameterized based on either ion- or salt-specific parameters, and this flexibility can greatly simplify the thermodynamic modeling of mixtures with limited number of ions. A
simple temperature dependence of the ion-solvent interaction
parameter enables prediction of thermodynamic properties of
complex mixtures over wide temperature and pressure
ranges.



\cite{Sun2019}: DTU eCPA salt-specific. It can be concluded that e-CPA gives
accurate agreements with the experimental solubility data of
H 2 O-gas and H 2 O-salt binary systems. With temperature-dependent ion-gas interaction parameters, eCPA gives
excellent agreement with the experimental gas solubility data
over wide ranges of temperature, pressure, and salt morality.  T < 523 K AAD = 6.6 \%.


- \cite{Li2020}: A promising flash technique at given moles, volume and temperature, known as NVT flash, is employed and the salting-out effect is reproduced by correcting the chemical potential of aqueous nonelectrolyte components.

-Another promising geological site for CO 2 sequestra-
tion is deep saline aquifers, which provide substantial storage ca-
pacity due to its large pore volume and wide distribution [2,9,10].
Most of the injected CO 2 is trapped in saline water by dissolution
and such a mechanism is called solubility trapping [11,12]. How-
ever, saline water usually has a considerable salt content and the
presence of salts could significantly reduce CO 2 solubility, which is
known as the salting-out effect. slug [20]. Moreover, unlike other gases, the
dissolution of CO 2 into aqueous phase often causes a density in-
crease, which can induce natural convection and facilitate CO 2
mixing with water [10,21,22]. There is also evidence that CO 2 can
form weak hydrogen bonds in the presence of associating species
[40].

-The salting-out effect is reproduced by
introducing the Debye-Hückel electrostatic contribution to chem-
ical potential of nonelectrolyte components in the aqueous phase.

- More like a correlation: uses the DH activity coefficient to modify the chemical potential of nonelectrolyte components. Mixture parameters (p): kcw (2 p), Scw (3 p), Uws (5 p), Ucs (3 p). Non-primitive model for the dielectric constant.


\cite{Springer2012}: - Usa gamma-phi approach. Within the Mixed-Solvent Electrolyte
(MSE) model framework the standard-state properties are
calculated from the Helgeson–Kirkham–Flowers equation of state whereas the excess Gibbs energy includes a long-range electrostatic interaction term expressed by a Pitzer–Debye–Hückel equation, a virial coefficient-
type term for interactions between ions and a short-range term for interactions involving neutral molecules. A large set of fitting parameters is required.

- data for H2O concentration in CO2-rich phase (CB mixture)



\cite{Bian2019}: The model, i.e., the thermodynamic model, focused on
the cubic-plus-association equation of state combined with the Wong−
Sandler mixing rule with six parameters is presented (CPA-WS) with composition dependency on kij. The model confirmed that the inert association scheme is more suitable for carbon dioxide (CO 2 ) than the 3B association scheme, for the level equilibrium in the carbon dioxide−water (H 2 O) system. CO 2 , as a non-self-associating compound (inert), only accepts a pair of electrons to solvate in water. It is necessary to consider cross-association with H 2 O. Meanwhile, the CO 2 −H 2 O−sodium chloride (NaCl) system can be reasonably well predicted by using two adjustable parameters of the H 2 O−NaCl and CO 2 −NaCl systems (DOES NOT ACCOUNT FOR ANY ELECTROSTATIC INTERACTIONS IN THE ELECTROLYTE SOLUTION).


\cite{Chabab2019}:For the modeling part, an electrolyte version of the PR-CPA (Peng-Robinson Cubic Plus Association) equation of state was developed under the name “e-PR-CPA” (with MSA for ion-ion), as well as modified Soreide and Whitson (m-SW) model was used and improved. These two models using the phi-phi approach are compared with two geochemical models (Corvisier 2013 and Duan el al. 2006) using the gamma-phi approach, and against the literature and measured data. For the H 2 O-CO 2 binary system, solvation was considered by assuming that CO 2 has only one electron acceptor (ea) site (type 0ed-1ea) and that H 2 O has 4 association sites (type 4C) with two electron donor (ed) sites and two electron acceptor sites.

- In this work, we neglect the formation of ion
pairs, that occur at very high temperatures (above 300 °C) (Anderko
and Pitzer, 1993).

- apresenta otimos results porem, USA um kij para fase aquosa e outro pra fase CO2-rich, além disso USA dependencia de ms no paramtro a da agua



\cite{Spycher2010}: gama-phi: CORRELATIONS are presented to compute the mutual solubilities of CO 2 and chloride brines at temperatures 12–300 ◦ C, pressures 1–600 bar (0.1–60 MPa), and salinities 0–6 m NaCl. The formulation is computationally efficient and primarily intended for numerical simulations of CO 2 -water flow in carbon sequestration and geothermal studies. The phase-partitioning model relies on experimental data from literature for phase partitioning between CO 2 and NaCl brines, and extends the previously published correlations to higher temperatures. The model relies on activity coefficients for the H 2 O-rich (aqueous) phase and fugacity coefficients for the CO 2 -rich phase. Activity coefficients are treated using a Margules expression for CO 2 in pure water, and a Pitzer expression for salting-out effects. Fugacity coefficients are computed using a modified Redlich–Kwong equation of state and mixing rules that incorporate asymmetric binary interaction parameters (kij depends on concentration).

EGS: An example of multiphase flow simulation implementing the mutual solubility model is presented for the case of a hypothetical, enhanced geothermal system where CO 2 is used as the heat extraction fluid. In this simulation, dry supercritical CO 2 at 20 ◦ C is injected into a 200 ◦ C hot-water reservoir. Results show that the injected CO 2 displaces the formation water relatively quickly, but that the produced CO 2 contains significant water for long periods of time. The amount of water in the CO 2 could have implications for reactivity with reservoir rocks and engineered materials.

EGS: Rising costs of fossil fuels have also led to a renewed interest in
geothermal energy, and particularly in enhanced geothermal systems (EGS) (MIT 2006). With these systems, the recovery of heat from the subsurface is enhanced by water injection e.g., Goyal 1995; Stark et al. 2005) and in some cases, water is injected, circulated, and recovered from areas that originally are dry or contain very little water (e.g., MIT 2006). The
use of CO 2 instead of water as the heat transmission fluid in EGS has been proposed as a means not only to produce “clean” energy, but also to potentially sequester CO 2 through fluid losses at depth, which typically occur in any type of EGS (Brown 2000; Fouillac et al. 2004;
Pruess 2006, 2008). Other potential advantages of CO 2 -based EGS include increased heat extraction rates and wellbore flow compared to water-based systems, and lesser potential for unfavorable rock–fluid chemical interactions in engineered systems (Pruess 2006, 2008). The solubility of CO 2 in water has obvious implications for long-term carbon sequestra-
tion and water–rock interactions (e.g., Gilfillan 2008; Moore et al. 2005; Rosenbauer et al. 2005; Palandri et al. 2005; Kaszuba et al. 2003). The solubility of water into CO 2 is also of importance, because it affects the reactivity of CO 2 with surrounding rocks (e.g., Regnault et al. 2005; Rimmelé et al. 2008; Lin et al. 2008; Suto et al. 2007) and determines the capacity of injected CO 2 to dry the rock formations adjacent to injection wells (Giorgis et al. 2007; Hurter et al. 2007; Pruess and Müller 2009). In the case of a CO 2 -based EGS, the partitioning of water into CO 2 would also determine the time required for the removal of water from the (central portion of the) EGS, and the type of chemical interactions that may take
place between the CO 2 plume, reservoir rocks, and engineered systems. Operating enhanced geothermal systems (EGS) with CO 2 instead of water as heat transmission fluid is a novel concept that was originally proposed by Brown (2000).

- Simulação de EGS em um grid 2d: These results suggest that although water is relatively quickly displaced by the CO 2 plume, significant amounts of water may remain dissolved in the produced CO 2 phase for long periods of time. The presence of water in the CO 2 would have implications not only
for the design of heat extraction systems, but would also dictate the amount of reactivity of CO 2 with the reservoir and engineered systems. For these reasons, future assessment of the potential for CO 2 -EGS will need to focus not only on flow/recovery issues, but also on the reactivity of CO 2 containing various amounts of water (including dry CO 2 ) with reservoir
rocks and other relevant materials.




\cite{Sun2021}:CORRELATION WITHOUT ANY PHYSICAL MEANING. Reliable and accurate prediction of phase-partitioning behaviors in CO 2 -Brine system over a wide range of temperature, pressure and ionic composition (T–P–m salts ) is of fundamental and practical importance to many geological-, energy-, and environmental-related technical applications.




\cite{Xiong2023}: MUITO PARECIDO COM MINHA ECPA (COM KIJ E SIJ, T < 473 K):A modified electrolyte cubic plus association equation of state (eCPA EoS) is proposed here as a suitable model for CO 2 Capture and Storage (CCS) fluids to model CO 2 , CH 4 , H 2 S, N 2 , O 2 , Ar, and air solubility in multi-component brines containing NaCl, KCl, CaCl 2 , MgCl 2 , and Na 2 SO 4 . A new correlation of the g- factor of water as a function of temperature is presented to accurately calculate the dielectric properties of pure water with CPA EoS and a new formula related to ions associated with polar molecules is presented to describe the mixed solvents containing salts by neglecting ion–solvent association, which avoid the computational complexity of hydrogen bond formation. The performance of the modified eCPA model is improved for modeling the gas–water system by considering the cross association between H 2 S/CO 2 and water. The modified eCPA model can accurately calculate the gas solubility in brine and water content in gas phase for water-salt-gas systems.

- Maribo-Mogensen (Maribo-Mogensen et al., 2013) calculated the fraction of component i that is not bound to an ion by assuming that ions associate with polar molecules. However, the ion–solvent association is very complex. Hence we neglect the ion–solvent association and derive a simpler formula (no need to compute the chi derivatives). we neglect ion-solvent association by
assuming the association strength is the same as the molar volume (???)

- The performance of eCPA model is improved for modeling gas–
water systems by considering the cross-association between
H 2 S (or CO 2 ) and H 2 O. The presented eCPA model can accurately predict the gas solubility in brine and water content in gas phase for water-salt-gas
systems. The presented eCPA EoS can accurately predict gas solubility in
multi-salt systems.



\cite{Yang2024}: A new association model, CPA-MHV1, is developed by combining the Soave-Redlich-Kwong (SRK) equation of state (EoS) with the Cubic-Plus-Association (CPA) EoS based on Michelsen’s improved Huron-Vidal mixing rule. . The results demonstrate that the best performance is achieved when solvation between CO 2 and H 2 O is considered. There is two non-symmetric binary adjustable parameters, τ ij and τ ji and Sij.

Usa TAK e TOD quando melhor convem.

- Does not account for any particular electrostatic interactions when there is salt in the solution (treat as neutral particles). T < 543 K (CW), T < 473 K (CB)






















MONTE CARLO SIMULATIONS

\cite{Vorholz2000}: NVT- and NpT-Gibbs ensemble Monte Carlo Ž GEMC . simulations were applied to describe the vapor–liquid equilibrium of water Ž between 323 and 573 K . , carbon dioxide Ž between 230 and 290 K . and their binary mixtures Ž between 348 and 393 K . For water, the
SPC- and TIP4P models give superior results for the vapor pressure when compared to the SPCrE model. The
vapor–liquid equilibrium of the binary mixture, carbon dioxide–water, was predicted using the SPC- as well as
the TIP4P model for water and the EPM2 model for carbon dioxide.

The enthalpy of vaporization as calculated from the different water models is plotted in Fig. 7. The SPC/E model systematically overestimates the experimental value over the whole simulated temperature range.


\cite{Vorholz2004}: - The solubility of carbon dioxide in aqueous solutions of sodium chloride is studied by NVT- and NpT-Gibbs Ensemble Monte Carlo Simulations at 373, 393 and 433 K at pressures up to 10 MPa. The intermolecular forces are approximated by effective pair potentials (SPC and TIP4P models for water, the EPM2 potential model for carbon dioxide and several pair potentials for sodium chloride). Unlike interactions are estimated applying common combining rules. The simulation results are compared with experimental data. The experimentally observed “salting-out effect” is predicted by the simulation, but the decrease of the solubility of carbon dioxide caused by the presence of salt is overestimated. An improvement is achieved by introducing binary interaction parameters into unlike pair potentials. The SPC/E model was not employed in the present simulations since its vapor pressure predictions are significantly too low (40–80 \%).

-Equilibrium was assumed to be reached, when the instantaneous thermodynamic properties were fluctuating around constant values. The statistical uncertainty of the simulation results was estimated by dividing the whole production period into up to 25 (depending on the length of the production period) blocks of equal size and calculating the standard deviation of the block averages. Since it is experimentally known that the concentration of sodium chloride in the gas phase can be neglected, the transfer probability for ions was set to zero in the simulations. Although the transfer probability of sodium chloride is set to zero its molality in the liquid phase does not remain constant. This is caused by the transfer of water molecules from/to the aqueous phase.

- The simulation results clearly reveal a salting-out effect. However, the predictions overestimate the decrease of the liquid phase concentration of carbon dioxide. For example, at 393 K and mNaCl = 4.6 mol/kg the decrease is overestimated by 30–50 \%, at mNaCl = 6.8 mol/kg the difference is between 43 and 58 \%. The statistical uncertainties of the simulation results for the carbon dioxide solubility in the liquid phase are smaller, i.e. between 27 and 46 \%. The influence of sodium chloride on the solubility is overestimated. This might also be caused by the fact that the influence of polarizability is neglected in the pair potential models. The agreement with experimental data could be improved by introducing a binary interaction parameter for the Lennard-Jones cross diameters of carbon dioxide and sodium chloride.



\cite{Liu2011}: Histogram-reweighting grand canonical Monte Carlo simulations were used to obtain the phase behavior of CO 2 -H 2 O mixtures over a broad temperature and pressure range (50 C e T e 350 C, 0 e P e 1000 bar). We performed a comprehensive test of several existing water (SPC, TIP4P, TIP4P2005, and exponential-6) and carbon dioxide (EPM2, TraPPE, and exponential-6)
models using conventional LorentzBerthelot combining rules for the unlike-pair parameters. None of the models we studied reproduce adequately experimental data over the entire temperature and pressure range, but critical assessments were made on the range of T and P where particular model pairs perform better. Near the critical region (250 C < T e 350 C), the critical points are not accurately estimated by any of the models studied, but the exponential-6 models are able to qualitatively capture the critical loci and the shape of the phase envelopes. Local improvements can be achieved at specific temperatures by introducing modification factors to the Lorentz-Berthelot combining rules, but the modified combining rule is still not able to achieve global improvements over the entire temperature and pressure range. 

- Using histogram-re-weighting Grand Canonical Monte Carlo (GCMC) techniques 
and mixed-field finite size scaling analysis, 27 we are also able to obtain the critical points at high temperatures.

- The gas-phase data are more difficult to assess, due to the lack
of agreement between experimental data sets from different
laboratories. 

- The TraPPE/TIP4P2005 and the EPM2/TIP4P2005 models pro-
vide better prediction of the liquid branches at 350 and 300 C
but not of the gas phase data, and they are both incapable of
predicting the intrinsic critical points at 300 C and below. On
the other hand, the EPM2/TIP4P and EPM2/SPC models do
not exhibit phase coexistence above T = 275 C. At T = 275 C,
they both have critical pressures that are much lower than the
experimental values. One of the reasons for most models to yield poor predictions is their generally poor agreement with experimental vapor
pressures. The critical points of the exponential-6 and TIP4P2005 models are very close to the experimental value, which could explain the fact that the
binary phase diagrams using these models better resemble the
experimental results at high temperatures(e.g., T = 350 C)

- There may be system size effects if the system is lower than 8sigma-CO2, higher than that there is no noticeable effect

- All of the models we tested have simple fixed-charge interac-
tions, effective pair potentials and commonly adopted mixing
rules, all of which could limit predictive capability. Further im-
provement of the quality of predictions using molecular simula-
tions may result from testing the effect of polarizability, many-
body interactions, and different combining rules in predicting the
phase behavior of pure systems and mixtures. 



\cite{Liu2013}: Direct interfacial molecular dynamics simulations are used to obtain the phase behavior and interfacial tension of CO 2 –H 2 O–NaCl mixtures over a broad temperature and pressure range (50 C < T < 250 C, 0 < P < 600 bar) and NaCl concentrations (1–4 mol/kg H 2 O). The predictive ability of several existing water (SPC and TIP4P2005), carbon dioxide (EPM2 and TraPPE), and sodium chloride (SD and DRVH) models is studied and compared, using conventional Lorentz–Berthelot combining rules for the unlike-pair parameters. Under conditions of moderate NaCl molality (1 mol/kg
H 2 O), the predictions of the CO 2 solubility in the water-rich and CO 2 -rich phase resemble those in the CO 2 –H 2 O system. Consistent with our previous work, the TraPPE/TIP4P2005 model combination gives the best overall performance in predicting coexistence composition and pressure in the water-rich phase. Critical assessments are also made on the ranges of temperature and pressure where particular model combinations work better.

- Histogram-reweighting techniques provide a more accurate route to calculate the phase coexistence properties near critical points but are not suitable for the study of solutions with added salt, due to the extremely
low acceptance ratio of adding or deleting strongly interact-
ing electrolytes. In recent years, direct interfacial simulation techniques have become a robust and simple tool to obtain phase equilibrium properties
of both pure fluids 28–30 and mixtures, 31–35 primarily because
computational capabilities are now adequate for simulating
relatively large systems over the appreciable times needed to
reach equilibration through diffusion. Another important property of the vapor–liquid interface is the surface tension. This quantity can be readily evaluated using its mechanical definition involving knowledge of the
diagonal elements of the pressure tensor

- Gibbs ensemble Monte Carlo simulations applied to such systems often
assume that no salt is present in the vapor phase. However,
such an assumption will not hold near the critical locus
where the compositions of the two phases become indistin-
guishable. Direct interfacial MD simulations provide a fast
and simple way to check the range of applicability of Gibbs
ensemble simulations. Our work also points out the limitations of the existing models with fixed-charged, additive pair interactions in pre-
dicting phase equilibrium and interfacial properties of com-
plex CO 2 –H 2 O–NaCl mixtures. We also point out that the
currently widely adopted Lorentz–Berthelot combining rule
is not sufficient to describe the contributions of mixing of
unlike molecules to the interfacial properties.



\cite{Vlcek2011}: NPT-GEMC, T = 29, 323, 348 K e P = 100, 200, 400 bar. The unlike-pair interaction parameters for the SPC/EEPM2 models have been optimized to reproduce the mutual solubility of
water and carbon dioxide at the conditions of liquidsupercritical fluid phase equilibria. An efficient global optimization of the
parameters is achieved through an implementation of the coupling parameter approach, adapted to phase equilibria calculations in
the Gibbs ensemble, that explicitly corrects for the overpolarization of the SPC/E water molecule in the nonpolar CO 2
environments. The resulting H 2 O-CO 2 force field accurately reproduces the available experimental solubilities at the two fluid
phases in equilibria as well as the corresponding species tracer diffusion coefficients.

- The need for an adequate description of the unlike-pair
interactions between different molecular species in a fluid
mixture comes from the fact that those interactions are the actual
sources of thermodynamic nonideality.

- The SPC/E model, in contrast to its predecessor the SPC
model, involves an additional “polarization correction” to the
heat of vaporization to account for the “missing” water self-
polarization in the liquid phase, which results in an additional
enhancement of the permanent dipole moment from 2.27D
(SPC) to 2.35D. While this modification greatly improves
various liquid properties (i.e., the original intended target), it
also causes significant deviations from the experimental values of
the orthobaric vapor densities and corresponding equilibrium
pressure because the fixed-charge model (with its resulting
augmented dipole moment) represents an unrealistic overpolar-
ized environment. Obviously, water in either a
pure vapor phase or as a highly dilute solute in a nonpolar dense
solvating environment exhibits comparatively smaller polariza-
tion, and consequently, a fixed-charge water model whose
parameters were adjusted to describe liquid-like aqueous envir-
onments must be “depolarized” in order to properly describe real
water in a nonpolar environment. The overpolarization has a significant impact on the predicted water solubility in less polar environments such as those involving single phases of either carbon dioxide in a water-rich phase
or water in a carbon dioxide-rich phase, and two (vaporliquid
or liquidliquid) phases in equilibrium. Fo

- If a hydrogen bond is geometrically defined as an O-H pair at separa-
tions less than 2.5 Å, 56 then a dissolved water molecule in the
CO 2 surroundings would form 1.1 bonds (at 298 K and 400
bar). In contrast, in the water-rich phase the hydrogen bonding
capacity of a water molecule can be saturated by other water
molecules, resulting in less opportunity for CO 2 oxygens to
participate in such interactions. On average, a dissolved CO 2
molecule forms 0.7 hydrogen bonds.

-We found that when the self-polarization correction for the chemical potential of SPC/E water in the CO 2 -rich phase is included, the predicted
solubilities are generally within the simulation/experimental
combined uncertainties.



\cite{Orozco2014}: - Monte Carlo simulations in the Gibbs ensemble
were used to obtain optimized intermolecular potential
parameters to describe the phase behavior of the mixture
CO 2 /H 2 O, over a range of temperatures (T = 423, 473, 523 K) and pressures  (P = 200, 400, 600, 800 bar) relevant
for carbon capture and sequestration processes. Commonly used
fixed-point-charge force fields that include Lennard-Jones 12−6
(LJ) or exponential-6 (Exp-6) terms were used to describe CO 2
and H 2 O intermolecular interactions. For force fields based on
the LJ functional form, changes of the unlike interactions
produced higher variations in the H 2 O-rich phase than in the
CO 2 -rich phase. A major finding of the present study is that for
these potentials, no combination of unlike interaction
parameters is able to adequately represent properties of both
phases. Changes to the partial charges of H 2 O were found to produce significant variations in both phases and are able to fit
experimental data in both phases, at the cost of inaccuracies for the pure H 2 O properties. By contrast, for the Exp-6 case,
optimization of a single parameter, the oxygen−oxygen unlike-pair interaction, was found sufficient to give accurate predictions
of the solubilities in both phases while preserving accuracy in the pure component properties. These models are thus
recommended for future molecular simulation studies of CO 2 /H 2 O mixtures.

- The TraPPE −SPC/E (LJ) model slightly underestimates the solubilities in
the aqueous phase, and significantly underestimates the
solubilities in the CO 2 phase, over the whole range of
pressures and temperatures studied. 

- Since the inclusion of flexibility is a different way to change the dipole moment in the system, without directly modifying the original model parameters, we considered important to study such models. we have found that instead of improving the solubility predictions the flexibility effect produces solubilities even lower than those obtained using the rigid molecule. This is expected, since according to the results presented previously, a lower dipole moment is needed to increase the solubilities in
both phases.

- One of the disadvantages of the models used in this study
is related to the fact that polarizability is not explicitly
considered.



\cite{Jiang2017a}- NPT-GEMC T = 423, 473, 523 K) and pressures  (P = 200, 400, 600, 800 bar). Phase equilibria of water/CO 2 and water/n-alkane
mixtures over a range of temperatures and pressures were
obtained from Monte Carlo simulations in the Gibbs ensemble.
Three sets of Drude-type polarizable models for water, namely the
BK3, GCP, and HBP models, were combined with a polarizable
Gaussian charge CO 2 (PGC) model to represent the water/CO 2
mixture. The HBP water model describes hydrogen bonds
between water and CO 2 explicitly. All models underestimate
CO 2 solubility in water if standard combining rules are used for
the dispersion interactions between water and CO 2 . With the
dispersion parameters optimized to phase compositions, the BK3
and GCP models were able to represent the CO 2 solubility in
water, however, the water composition in CO 2 -rich phase is
systematically underestimated. Accurate representation of compo-
sitions for both water- and CO 2 -rich phases cannot be achieved even after optimizing the cross interaction parameters. By
contrast, accurate compositions for both water- and CO 2 -rich phases were obtained with hydrogen bonding parameters
determined from the second virial coefficient for water/CO 2

- For phase equilibria at elevated temperatures (>473 K) and pressures (>200 bar), which are of particular interest for geochemical applications,
there are mainly three sets of experimental measurements from
Takenouchi and Kennedy, 2 Todheide and Frank, 3 and
Capobianco et al. 4 Although these show consistent results for
CO 2 solubility in water, the reported equilibrium compositions
of the CO 2 -rich phase are in disagreement.

- Duan et al. proposed a series of important
correlations 6−8 to calculate CO 2 solubility in water as well as
in aqueous NaCl solutions up to 2000 bar. Equations of state
have also been developed for the water/CO 2 mixture. Models
rooted in statistical mechanics, such as statistical associating
fluid theory (SAFT) 9−15 and cubic plus association (CPA)
equations of state, 16−18 are better suited for description of this
system than purely empirical models, because of its highly
nonideal character. Although these equations of state show
remarkable accuracy in calculation or prediction of phase
equilibria for the water/CO 2 mixture, temperature- or pressure-
dependent model parameters are required, and these
parameters need to be obtained by fitting to mixture
experimental data. Molecular modeling and simulation using transferrable
molecular models can provide useful information for thermophysical properties of the water/CO 2 mixture.

- the Exp-6 water model is inaccurate for the structure of liquid
water, and it gives poor performance with respect to the
prediction of thermodynamic properties for the water+NaCl
mixture, 31 which is also important to CO 2 geological storage
and other geochemical applications.

- The dielectric permittivities of water- and CO 2 -rich phases
differ significantly, but the nonpolarizable models are unable to
respond the change of the system electric environment, as they
carry a fixed dipole moment. Therefore, the failure of fixed-
point-charge literature models may be attributed to the lack of
polarizability. 29,30 However, another important intermolecular
interaction that is missing in most of the existing models is the
charge transfer or hydrogen bonding interaction between water
and CO 2. O 2 . For the HBP water model, a short-range and orientation-dependent hydrogen bonding term is included to account for the effect of charge transfer during the formation of hydrogen bond,



\cite{Blasquez2023}: In this work, we performed Molecular Dynamics (T = 323 K, P = 200 bar, ms = 1, 3 m) simulations to investigate the solubility of CO 2 in both pure water and NaCl aqueous solutions, employing various force fields for water and NaCl. We first focussed on the solubility of CO 2 in TIP4P/2005 and TIP4P/Ice water models. We found that it was necessary to apply positive deviations to the energetic Lorentz-Berthelot combining rules for the cross interactions of water and CO 2 in both models to accurately replicate the experimental solubility of the gas. Then we found a negative adsorption of ions at the NaCl solution–CO 2 interface. We also observed an increase of the interfacial tensions of the system when adding salt which is related with the observed negative adsorption of the ions at the interface. Furthermore, we found that unit charge models tend to highly overestimate the change in the interfacial tension compared to scaled charge models. Finally, we also explored the salting out effect of CO 2 using different force fields for NaCl. Our findings indicate that unit charge force fields significantly overestimate the salting out effect of CO 2 while the scaled charge model Madrid-2019 accurately reproduces the experimental salting out effect of CO 2 in NaCl solutions.

- In the case of the salting out effect, the presence of
ions weakens the interaction between the solute and the
water molecules (in the sense of modifying the solva-
tion free energy of the solute to make it less favourable),
leading to a decreased solubility of gases in the solu-
tion.

- Joung and Cheatham (JC) model [41] or the Smith and Dang (SD)
model [42]. Indeed, these models fall short in reproduc-
ing crucial transport properties like viscosity, diffusion
coefficient, or electrical conductivity [43–46]. Addition-
ally, these models tend to underestimate salt solubilities
when compared to experimental values [44, 47].

- the energy parameter was obtained by multiplying by χ =
1.11 the value obtained using default Lorentz-Berthelot
energetic combining rules acc